\chapter{桌面：Hyprland 与 Niri}
\chapterlead{图形界面是 Linux 之上的一层普通软件。本章用两个平铺式合成器搭建桌面，说明 Wayland 下各部件如何分工，以及中文输入与显示配置的位置。GNOME 与 KDE Plasma 作为对比出现。}

\section{起步：最小可用的 Hyprland}

Hyprland 是 Wayland 合成器，负责绘制窗口、处理输入、管理布局。平铺指窗口自动占据屏幕，不互相重叠。

\begin{nixcode}
programs.hyprland.enable = true;

services.displayManager.ly.enable = true;

environment.sessionVariables = {
  NIXOS_OZONE_WL = "1";
  XDG_SESSION_TYPE = "wayland";
};
\end{nixcode}

重建后从显示管理器登录，选择 Hyprland 会话。首次进入是空白桌面，按 \texttt{Super+Return} 打开终端（默认绑定需要终端已安装）。

配置写在 \texttt{\textasciitilde/.config/hypr/hyprland.conf}，用 Home Manager 管理更稳：

\begin{nixcode}
wayland.windowManager.hyprland = {
  enable = true;
  settings = {
    monitor = [ ",preferred,auto,1" ];
    "$mod" = "SUPER";
    bind = [
      "$mod, RETURN, exec, kitty"
      "$mod, Q, killactive"
      "$mod, E, exec, thunar"
      "$mod, 1, workspace, 1"
    ];
    general = {
      gaps_in = 4;
      gaps_out = 8;
      border_size = 2;
    };
  };
};
\end{nixcode}

\texttt{bind} 的四个字段是：修饰键、按键、动作、参数。改完执行 \texttt{home-manager switch} 或 \texttt{hyprctl reload} 生效。

\section{Niri：另一种平铺思路}

Niri 同样平铺，但采用可滚动的列布局：窗口按列排列，屏幕放不下时整行横向滚动。它适合窗口数量多、经常在几件事之间切换的场景。

\begin{nixcode}
programs.niri.enable = true;
\end{nixcode}

\begin{textcode}
// ~/.config/niri/config.kdl
input {
    keyboard {
        xkb { layout "us,cn"; }
    }
}

binds {
    Mod+Return { spawn "kitty"; }
    Mod+Q { close-window; }
    Mod+Left  { focus-column-left; }
    Mod+Right { focus-column-right; }
}
\end{textcode}

Niri 的配置用 KDL 格式，与 Hyprland 的写法不同，但概念对应：修饰键加按键绑定动作。两者可以装在同一台机器上，登录时选择会话。

\section{Wayland 的部件分工}

\begin{booktable}{\textwidth}{J{0.846} J{1.154}}
\toprule
部件 & 职责 \\
\midrule
内核与驱动 & 显卡、输入设备 \\
合成器 & 绘制窗口、合成画面、处理输入（Hyprland、Niri） \\
会话与登录 & 启动合成器，管理用户会话（ly、greetd、GDM） \\
通知与状态栏 & 通知、音量、网络图标（mako、waybar） \\
输入法 & 中文输入（fcitx5、ibus） \\
\bottomrule
\end{booktable}

平铺合成器不自带状态栏与通知，需要另配。最小组合是 waybar 加 mako：

\begin{nixcode}
programs.waybar.enable = true;
services.mako.enable = true;

# 截图与剪贴板
environment.systemPackages = with pkgs; [
  grim slurp wl-clipboard
];
\end{nixcode}

\texttt{grim} 抓屏，\texttt{slurp} 选区域，两者组合成截图命令：\texttt{grim -g "\$(slurp)" - | wl-copy}。

\section{中文显示与输入}

字体、区域设置与输入法是三件事。

\begin{nixcode}
i18n.defaultLocale = "zh_CN.UTF-8";
i18n.inputMethod = {
  enable = true;
  type = "fcitx5";
  fcitx5.addons = with pkgs; [ fcitx5-chinese-addons ];
};

fonts.packages = with pkgs; [
  noto-fonts-cjk-sans
  noto-fonts-cjk-serif
  sarasa-gothic
];
\end{nixcode}

\texttt{defaultLocale} 决定系统语言。输入法由 fcitx5 提供，装好后在合成器里设置环境变量 \texttt{GTK\_IM\_MODULE=fcitx}、\texttt{QT\_IM\_MODULE=fcitx}，否则部分程序收不到输入法事件。中文显示依赖字体，缺字时表现为方框，装字体即可。

\section{GNOME 与 KDE Plasma}

两者是完整的桌面环境，自带窗口管理、状态栏、设置面板与一套应用。配置成本低，代价是可定制性弱于平铺合成器。

\begin{nixcode}
services.desktopManager.gnome.enable = true;
# 或
services.desktopManager.plasma6.enable = true;
\end{nixcode}

选择建议：日常以浏览器和 IDE 为主，GNOME 够用；需要大量窗口切换与键盘操作，Hyprland 或 Niri 更省事；需要图形化设置每一处细节，KDE Plasma 最完整。三者可以并存，登录时切换，命令行部分不受影响。

\section{练习}

\begin{enumerate}
  \item 启用 Hyprland 并登录，用 \texttt{Super+Return} 打开终端，记录默认配置文件的路径。
  \item 给 Hyprland 增加三个快捷键：打开文件管理器、切换工作区、退出当前窗口。
  \item 安装 Niri，比较同一操作在两个合成器下的默认键位差异。
  \item 配置 fcitx5，在终端与浏览器里各输入一段中文，确认都能上屏。
  \item 用 \texttt{grim -g "\$(slurp)" - | wl-copy} 截图并粘贴，说明管道中每一步的作用。
\end{enumerate}

\begin{checkpoint}
本章结束时，应能：
\begin{itemize}
  \item 说明 Wayland 下合成器、登录管理、状态栏、输入法各自做什么；
  \item 写出 Hyprland 与 Niri 的最小可用配置；
  \item 配置中文区域、字体与输入法，并解释三者缺一的表现；
  \item 列出平铺合成器相对完整桌面环境的收益与代价；
  \item 判断自己应当选哪一种桌面。
\end{itemize}
\end{checkpoint}
